% =====================================================================
%  CH 12-16 : Gouvernance · Menaces et limites · Jalons · Correspondance · Décisions
% =====================================================================
\section{Gouvernance}\label{sec:gouvernance}

\subsection{Journal des décisions — format normalisé}\label{sec:journal}

Le protocole ne se révisant jamais en silence, chaque décision structurante est consignée
dans un journal. La v0.4.0 normalise le format (D24) :

\begin{table}[htbp]
\centering\small
\caption{Format normalisé d'une entrée du journal des décisions.}
\label{tab:formatdecisions}
\begin{tabularx}{\textwidth}{@{}l Y@{}}
\toprule
\textbf{Champ} & \textbf{Exigence} \\
\midrule
ID & \code{DNNN}, séquentiel, jamais réutilisé \\
Date & date de consignation \\
Statut & active · amendée par Dxx · invalidée par Dxx \\
Contexte & exigence, incident ou évaluation qui motive la décision \\
Choix retenu & décision en une phrase, rattachée à la fonction objectif \\
Alternatives rejetées & au moins une, avec motif \\
Amendements & référence explicite à la décision remplacée, le cas échéant \\
Développement & argumentaire détaillé si la décision engage l'architecture \\
\bottomrule
\end{tabularx}
\end{table}

Trois règles : le journal est \textbf{append-only} (une entrée ne s'efface jamais) ; une
révision \textbf{invalide explicitement} ce qu'elle remplace ; aucune entrée n'est
rétrodatée. La conformité du journal — présence des champs, unicité des ID, chaînage des
amendements — est contrôlée mécaniquement par l'audit de gouvernance.

\subsection{Tableau de bord de conformité}\label{sec:tableaudebord}

La v0.4.0 unifie les indicateurs de santé en un \textbf{tableau de bord} (D24), généré par
la commande d'audit à chaque commit et à chaque cristallisation, archivé avec la carte :

\begin{table}[htbp]
\centering\small
\caption{Tableau de bord de conformité : métriques, cibles, sources.}
\label{tab:dashboard}
\begin{tabularx}{\textwidth}{@{}Y l Y@{}}
\toprule
\textbf{Métrique} & \textbf{Cible} & \textbf{Source de mesure} \\
\midrule
Blocs avec spécification I à jour & 100 \% & audit de traçabilité \\
Chaînes de traçabilité complètes I $\leftrightarrow$ C $\leftrightarrow$ B
$\leftrightarrow$ tests & 100 \% & audit de traçabilité \\
Construction & 0 / 0 & build \\
Tests différentiels & 100 \% verts & oracle \\
Empreintes des données embarquées & conformes & hachage vs carte \\
Défauts ouverts par classe T1–T6 & 0 & journal des décisions \\
Optimisations sans mesure (R11) & 0 & journal d'optimisation \\
Budgets tokens (localiser / éditer / ajouter) & $\leq$ 40 / 3 000 / 6 000 &
instrumentation (dès M1) \\
Suite de conformité par profil actif & verte & suite du profil \\
Changements hors journal des décisions & 0 & diff vs journal \\
Pré-versions correctement balisées (dev / rc / final) & 100 \% & registre des balises \\
\bottomrule
\end{tabularx}
\end{table}

Le tableau de bord est la contrepartie mesurable de la gouvernance : chaque affirmation du
protocole s'y retrouve sous forme de métrique avec sa cible et sa source — aucune ne
s'appuie sur une déclaration non vérifiable. Les métriques dépendantes de l'outillage
affichent leur statut « cible » jusqu'à leur premier mesurage.

\subsection{Continuité et anti-enfermement}

L'export de continuité (D16) demeure l'unique forme de transfert de l'état de
développement : conteneur chiffré et signé regroupant I, B, C et un manifeste
autodescriptif ; oracle régénéré depuis I à l'arrivée, jamais exporté ; import en cinq
étapes dont l'échec d'une seule est terminal. Sauvegarde 3-2-1 chiffrée et escrow de la
clé maîtresse inchangés.

\section{Menaces et limites}\label{sec:menaces}

Le principe directeur demeure : séparer l'obfuscation, qui protège la propriété
intellectuelle, du durcissement, qui protège l'exécution. Conformément au principe de
Kerckhoffs, la sécurité ne repose jamais sur le secret de la carte : elle survivrait à sa
divulgation.

\subsection{Modèle de menace (étendu en v0.4.0, D20)}

\begin{table}[htbp]
\centering\small
\caption{Modèle de menace v0.4.0 — les six ajouts en gras.}
\label{tab:menaces}
\begin{tabularx}{\textwidth}{@{}Y Y Y@{}}
\toprule
\textbf{Menace} & \textbf{Contre-mesure} & \textbf{Résiduel} \\
\midrule
Rétro-ingénierie de l'IP & carte privée ; binaire sémantiquement muet ; variantes par
build & temps acheté — s'érode avec les LLM de décompilation \\
Empoisonnement de carte & signature ; refus de carte non signée ; verrou croisé (R8, R9)
& détecté — session bloquée \\
Corruption de carte (accidentelle) & AEAD par fragment + signature + verrou & détecté au
chargement \\
Fuite de la carte & chiffrement ; sharding ; escrow & perte limitée à un fragment \\
Vol de l'export de continuité & chiffrement intégral + signature & indéchiffrable sans
la phrase \\
\textbf{Binaire altéré} & checksums par bloc + signature du manifeste au lancement
(R10 a, b) & détecté — arrêt avant tout effet \\
\textbf{Rollback non autorisé} & monotonie de version signée, plancher anti-rollback
(R10 c) & détecté — refus \\
\textbf{Agent défaillant} & R2 ; R5 ; oracle dérivé de I ; budgets & l'erreur ne devient
jamais permanente sans vérification \\
\textbf{Données embarquées modifiées} & empreintes dans la carte ; vérification au
lancement (R10 b) & détecté — arrêt propre \\
\textbf{Environnement non conforme} & contrôles de conformité au profil (R10 d) ;
fail-closed & refus propre avant exécution \\
Fuzzing boîte noire & contrats, oracle, tests par propriétés & indépendant du secret de
la carte \\
Redistribution frauduleuse & empreintes par build — canary trap & fuite traçable \\
Instrumentation dynamique & anti-debug léger par mesure de temps & assumé \\
\bottomrule
\end{tabularx}
\end{table}

Les six menaces en gras sont les additions de la v0.4.0 ; toutes sont neutralisées par des
mécanismes \textbf{locaux et vérifiables} — la contre-mesure anti-rollback ne dépend
d'aucun service en ligne (le chargement doit échouer proprement hors ligne), et les
vérifications de R10 s'exécutent avant toute entrée dans le chemin d'exécution.

\subsection{Limites assumées}\label{sec:limites}

La rétro-ingénierie est un objectif d'économie, pas d'impossibilité : le coût de l'analyse
d'un binaire strippé, sans structure, dont chaque exemplaire diffère, reste supérieur à sa
valeur — mais les modèles de langage appliqués à la décompilation progressent, et
l'avantage doit être réévalué à chaque génération. L'anti-debug reste léger et assumé.
La furtivité face aux EDR est hors périmètre : un outil légitime par construction, dont
la transparence vis-à-vis du système est totale. Enfin, la gouvernance du protocole a été
validée sur un cas unique (N = 1) — l'étude de cas RoomEq, dont les limites sont
documentées dans l'édition académique (D18) : la généralité reste à établir par
réplication, le cœur natif par les jalons M0 à M7.

\section{Critères d'acceptation des jalons M0 à M7}\label{sec:jalons}

Huit jalons, chacun avec des critères de sortie mesurables (détail dans
\code{PLAN\_MILESTONES.md}). \textbf{Conditions transverses} (v0.4.0) : aucun jalon n'est
franchi sans oracle vert, suite de conformité verte pour les éléments couverts, tableau de
bord conforme et carte signée.

\begin{table}[htbp]
\centering\small
\caption{Critères d'acceptation des jalons.}
\label{tab:jalons}
\begin{tabularx}{\textwidth}{@{}l Y@{}}
\toprule
\textbf{Jalon} & \textbf{Critères de sortie mesurables} \\
\midrule
M0 — chaîne artefact & exécution sur Windows 11 ; round-trip 100 \% ; carte v0 signée ;
exécutable $<$ 4 Kio ; R1–R3, R8–R9 dès le premier octet \\
M1 — blocs réutilisables & oracle 100 \% sur la mini-libc ; première cristallisation ;
adresses stables démontrées ; première mesure effective des budgets tokens \\
M2 — patch à chaud & remplacement du bloc d'écho en exécution, sans interruption ;
rollback démontré ; régression post-patch nulle \\
M3 — entrées non fiables & un million de requêtes fuzzées sans crash ; proxy
transparent ; contrats renforcés \\
M4 — tunnel chiffré & HTTPS en transit ; échec propre sur certificat invalide \\
M5 — outil complet & interception d'un flux navigateur en VM isolée ; patch à chaud du
journal en production ; empreintes par build actives \\
M6 — portabilité inter-OS & suite de conformité verte ; \textbf{diff nul sur la couche
I} ; parité inter-profils vérifiée ; même oracle vert sur les deux profils \\
M7 — portabilité inter-ISA & round-trip 100 \% sur la nouvelle ISA ; suite de conformité
verte ; patch à chaud conforme (scratch + \code{fence.i}) ; verrou multi-profils ; export
multi-profils vérifié \\
\bottomrule
\end{tabularx}
\end{table}

M0 est le jalon existentiel — il prouve la chaîne notation, octets, exécutable, exécution
sans compilateur. M2 est le jalon de vérité — le patch à chaud a motivé l'architecture.
M6 et M7 transforment la portabilité de promesse en propriété démontrée, en isolant chaque
variable ; la suite de conformité par profil en est désormais le critère d'acceptation
explicite.

\section{Correspondance v0.3 $\to$ v0.4.0}\label{sec:correspondance}

\subsection{Traçabilité des évolutions}\label{sec:evolutions}

\begin{table}[htbp]
\centering\small
\caption{Les dix-sept évolutions de la v0.4.0, tracées depuis la v0.3.}
\label{tab:correspondance}
\begin{tabularx}{\textwidth}{@{}l Y Y l@{}}
\toprule
\textbf{Évolution} & \textbf{En v0.3} & \textbf{En v0.4.0} & \textbf{Statut} \\
\midrule
Couche I normalisée & « contrats légers » sans gabarit & gabarit complet
(§\ref{sec:coucheI}) & renforcé (D19) \\
Taxonomie des défauts & absente & classes T1–T6 (§\ref{sec:taxonomie}) & nouveau (D19) \\
Threat model & 7 menaces & 13, dont 6 nouvelles & étendu (D20) \\
R10 & checksums au lancement & + signature, anti-rollback, conformité, fail-closed
explicite & renforcé (D20) \\
Optimisation prématurée & interdite implicitement & interdiction normative R11 &
explicité (D21) \\
Journal d'optimisation & pratique de démonstration & format OPT-NN & officialisé (D21) \\
Fragmentation & 4 fragments & + norme « aucune lecture totale » & consolidé (D24) \\
Localisation par intention & étape 2 implicite & balayage interdit, formalisé &
renforcé (D24) \\
Budgets tokens & 40 / 3 000 / 6 000 & inchangés, instrumentés & consolidé (D24) \\
Suite de conformité profil & implicite & suite obligatoire (§\ref{sec:conformite}) &
nouveau (D22) \\
Parité inter-profils & critères dispersés & définition formelle (§\ref{sec:parite}) &
formalisé (D22) \\
Reverse-specs & absente (pratique démo) & officialisées (§\ref{sec:reversespecs}) &
officialisé (D23) \\
Pré-versions & absente (pratique \code{-dev.1}) & dev / rc / final & officialisé (D23) \\
Fixtures partagées & absente (pratique installateur) & contractuelles & officialisé
(D23) \\
Traçabilité I–C–B–tests & implicite & table normative + audit & formalisé (D24) \\
Journal des décisions & pratique & format normalisé & normalisé (D24) \\
Tableau de bord & indicateurs dispersés & métriques unifiées & nouveau (D24) \\
\bottomrule
\end{tabularx}
\end{table}

\subsection{Modifications explicitement rejetées}\label{sec:rejets}

Conformément au cahier de révision, les modifications suivantes ont été examinées et
\textbf{rejetées, donc non appliquées} :

\begin{table}[htbp]
\centering\small
\caption{Modifications examinées et rejetées.}
\label{tab:rejets}
\begin{tabularx}{\textwidth}{@{}Y Y@{}}
\toprule
\textbf{Modification} & \textbf{Motif du rejet} \\
\midrule
Modifier la hiérarchie Sécurité $>$ Performance $>$ Optimisation $>$ Coût en tokens &
invariant de révision (D19) — condition même de l'arbitrage \\
Introduire une couche d'abstraction à l'exécution (bytecode, VM) & trahit la fonction
objectif Performance ; la parité s'obtient par vérification, pas par couche \\
Promettre la portabilité universelle immédiate & la portabilité se vérifie par profil
(M6/M7), elle ne se déclare pas \\
Introduire des mécanismes de furtivité face aux EDR & hors périmètre assumé —
transparence totale de l'outil \\
Affaiblir R5, R8, R9 ou R10 & la révision ne contient que des renforcements ; R1–R9
repris à l'identique \\
Anti-rollback par contrôle en ligne au chargement & R10 doit échouer proprement hors
ligne — aucune dépendance réseau \\
Optimisation validée sans mesure & R11 : hot path + mesure avant/après + contrat stable \\
\bottomrule
\end{tabularx}
\end{table}

\section{Journal des décisions}\label{sec:decisions}

Le journal complet vit dans \code{docs/DECISIONS.md} du dépôt du projet. Synthèse :
D1–D10 (v0.2) fixent l'oracle éphémère, la notation, les contrats légers, la
cryptographie, le verrou croisé, les variantes par build, l'interface système et la
granularité ; D11–D13 (infrastructure) la sauvegarde triple et systématique ; D14–D16
(v0.3) l'universalité par profils, la cryptographie paramétrable et l'export de
continuité ; D17 la régénération systématique du PDF formel ; D18 l'édition académique
(présentation IMRAD du noyau v0.3, sans valeur normative) ; D19–D24 (v0.4.0, format
normalisé) la présente révision :

\begin{itemize}
\item \textbf{D19} : couche I normalisée, taxonomie T1–T6, hiérarchie de la fonction
objectif verrouillée comme invariant de révision ;
\item \textbf{D20} : menace étendue (binaire altéré, carte corrompue, rollback,
agent défaillant, données embarquées, environnement) et R10 renforcée (checksums,
signature, anti-rollback, fail-closed explicite, vérification locale) ;
\item \textbf{D21} : R11 optimisation mesurée et journal d'optimisation OPT-NN ;
\item \textbf{D22} : suite de conformité par profil et parité inter-profils à contrat
constant ;
\item \textbf{D23} : officialisation des reverse-specs, des pré-versions dev / rc /
final et des fixtures partagées ; numérotation à trois segments ;
\item \textbf{D24} : traçabilité I $\leftrightarrow$ C $\leftrightarrow$ B
$\leftrightarrow$ tests, journal des décisions normalisé, tableau de bord de conformité.
\end{itemize}

Deux décisions postérieures à la révision complètent le tableau, sans valeur
normative : D25 confie la production du présent PDF formel à la chaîne
LaTeX/Tectonic (générateur ReportLab retiré, archives des versions antérieures
conservées) ; D26 consigne le jalon M0, première étape de la feuille de route,
menée à terme le 26 octobre 2026 (chaîne artefact implémentée, \code{hello.exe}
exécuté sur Windows 11, carte v0 signée, quatre critères de sortie acquis).

Une décision ne se modifie jamais en silence : la révision remplace l'entrée par une note
d'invalidation explicite référençant la décision qui la supplante. La gouvernance
s'applique à elle-même — la présente révision en est la démonstration.
